上述对 $f$ 的固有性假设可用 $f$ 与 $K$ 在无穷远处的正则性条件代替. 例如, 有下述结果（对 $K\in\operatorname{Ob}D(X_{\bar\eta},\Lambda)$）.

\subsubsection*{命题 2.1.9.}
\underline{假设 $f$ 有分解 $f=f_1k$}
\begin{center}
\begin{tikzcd}[column sep=huge]
X \arrow[r, hook, "k"] & X_1 \arrow[r, "f_1"] & S,
\end{tikzcd}
\end{center}
\underline{其中 $f_1$ 固有, $k$ 将 $X$ 等同于 $X_1$ 中一个相对正规交叉除子 $D$ 的补集}（因此, 按假设, $f_1$ 在 $D$ 的一个邻域上光滑）. \underline{若在 $D$ 的某个邻域 $U$ 上, 层 $\mathcal H^i(K)$ 在 $U-(D\cup X_0)$ 上局部常值, 且沿正规交叉除子 $D\cup X_0$ 温和分歧, 则 (2.1.8.1) 和 (2.1.8.2) 为同构.}

$K$ 的超上同调谱序列说明, 只须处理 $K$ 由位于零次的一个层 $F_\eta$ 组成的情形. 将 $S$ 换成它在某个有限扩张中的正规化（直接由定义可见这样做是允许的）, 再应用 Abhyankar 引理, 可假设 $F_\eta$ 是 $X$ 上某个层 $F$ 在 $X_\eta$ 上的限制, 该层在 $U\cap X$ 上局部常值, 且沿 $D$ 温和分歧. 使用 2.1.2 的记号, 设 $\overline F$ 为 $F$ 在 $\overline X$ 上的逆像, 并记 $\bar f:\overline X\longrightarrow\overline S$.

为证明 (2.1.8.1) 为同构, 只须证明
\[
R\bar f_*\bigl(R\bar j_*\bar j^*\overline F\bigr)
\]
与由 $\bar s\longrightarrow\overline S$ 给出的基变换交换. 在区别三角
\[
\longrightarrow\overline F\longrightarrow
R\bar j_*\bar j^*\overline F\longrightarrow\Delta\longrightarrow,
\]
中, 层 $\mathcal H^i(\Delta)$ 的支集包含于 $X_{\bar s}$; 对 $\Delta$ 而言, 所讨论的基变换陈述显然成立. 因此只须证明 $R\bar f_*(\overline F)$ 与由 $\bar s\longrightarrow\overline S$ 给出的基变换交换; 或者, 由固有态射 $f_1$ 的基变换定理, 只须证明 $Rk_*(\overline F)$ 与基变换 $X_{1\bar s}\longrightarrow X_1$ 交换. 这就是下面的引理 2.1.10.

为证明 (2.1.8.2) 为同构, 只须对层 $F$ 证明补零延拓 $k_!$ 与函子 $R\bar j_*$ 交换; 即在 $D$ 的一个邻域上有
\[
R\Phi(k_!F)=0\qquad\hbox{on }X_{1\bar s}.
\]
这就是下面的引理 2.1.11.

\subsubsection*{引理 2.1.10.}
\underline{设} $f:X\longrightarrow S$ \underline{光滑,} $D\subset X$ \underline{为相对正规交叉除子, 置} $U=X-D\mathrel{\mathop{\hookrightarrow}^{j}}X$, \underline{并设} $F$ \underline{为 $U$ 上局部常值且沿} $D$ \underline{温和分歧的 $\Lambda$-模层. 则} $Rj_*F$ \underline{的形成与任意基变换} $T\longrightarrow S$ \underline{交换.}

此问题在 $X$ 的平展拓扑下是局部的. 因此可假设存在态射 $g:X'\longrightarrow X$, 其中 $X'$ 为 $X$ 在一个沿 $D$ 温和分歧的有限平展满覆盖 $g':U'\longrightarrow U$ 中的正规化, 使 $g'^*F$ 在 $U'$ 上为常值. 由 Abhyankar 引理, 还可假设 $X'$ 在 $S$ 上光滑, 且 $U'$ 为 $X'$ 中某个相对正规交叉除子 $D'$ 的补集.

层 $F$ 单射地映入 $g'_*g'^*F$, 而余核仍满足引理的假设. 反复进行此构造（同一个 $g$ 始终适用）, 得到 $F$ 的一个消解, 其各项形如 $g'_*G$, 其中 $G$ 为常值且 $g'$ 如上所述. 因此只须对这样的层证明断言; 等价地, 只须对 $G$ 在 $U'$ 上为常值时的 $(X'/S,D',G)$ 证明断言. 这是 SGA~5, II, 4.4 的特例.

\subsubsection*{引理 2.1.11.}
\underline{设 $S$ 为 Hensel 离散赋值环谱, $f:X\longrightarrow S$ 光滑, $D\subset X$ 为相对正规交叉除子, 置 $U=X-D\mathrel{\mathop{\hookrightarrow}^{j}}X$, 并设 $F$ 为 $U$ 上局部常值且沿 $D$ 温和分歧的 $\Lambda$-模层. 则}
\[
\underline{R\Phi(j_!F)=0.}
\]

刚才使用的分解论证将我们归约到 $F$ 为常值层、值为 $G$ 的情形. 可假设 $D$ 是光滑除子 $(D_i)_{i\in I}$ 的并. 对 $J\subset I$, 以 $D_J$ 表示所有 $i\in J$ 的 $D_i$ 之交. 层 $F$ 有一个消解 $F^\bullet$, 其各项为层 $G_J$ 的和, 其中这些层在 $D_J$ 上为常值, 并补零延拓到 $X$. 因此只须对这样的一个层证明引理, 而这可由对 $D_J/S$ 应用 (2.1.5) 得到.

我猜想, 当 $S$ 和 $S'$ 为优良概形时, (2.1.7.5) 是同构. 下面给出一个部分结果（另见 2.4.2）.

\subsubsection*{命题 2.1.12.}
\underline{态射 (2.1.7.5) 在下列情形中为同构:}

\medskip
\noindent(a) $S$ 为优良概形, $S'$ 为其完备化;

\noindent(b) $S$ 和 $S'$ 均为等特征零.

在情形 (b) 中, 可归约到 $S$ 和 $S'$ 均严格局部, 且 $S$ 的一个一致化参数仍为 $S'$ 的一致化参数的情形（作标量扩张而不改变 $\bar S'\longrightarrow\bar S$）. 在 (a) 和 (b) 两种情形中, 若用上划线表示在分式域的可分闭包中的正规化, 则有
\[
\bar S'\simeq S'\times_S\bar S.
\]
回到定义, 只须证明 $S'\longrightarrow S$ 普遍局部无上同调（应用 SGA~4, XVI, 1.1）. 这由 I, 0.5.1（基于 Artin 版本的 N\'{e}ron 解消奇点）和 SGA~4, XV, 2.1（光滑态射的局部无上同调性定理）推出.

\subsubsection*{2.1.13.}
假设 $f:X\longrightarrow S$ 有限型. 由 2.1.4 和 SGA~4, XIV, 3.2（参见 I, 4.2(i)）, $R\Psi$ 的上同调维数（SGA~4, XVII, 1.2.9）至多为 $\dim(X_\eta)$. 因此 $R\Psi$ 可扩张到整个导出范畴,
\[
R\Psi:D(X,\Lambda)\longrightarrow D(X_s\times_sS,\Lambda),
\]
且将 $D^-$ 映入 $D^-$. 设 $L\in\operatorname{Ob}D^-(\Lambda)$ 为右 $\Lambda$-模的上有界复形, $K\in\operatorname{Ob}D^-(X,\Lambda)$. 则有同构
\begin{equation}\tag{2.1.13.1}
L\overset{\mathbf L}{\otimes}_\Lambda R\Psi(K)
\xrightarrow{\ \sim\ }
R\Psi\bigl(L\overset{\mathbf L}{\otimes}_\Lambda K\bigr).
\end{equation}
为验证 (2.1.13.1), 先将 $L$ 换成自由 $\Lambda$-模的复形. 因为 $R\Psi$ 是三角函子且上同调维数有限, 只须将 $L$ 换成其任一项后验证断言. 通过滤过余极限, 可归约到 $L$ 为有限自由模的情形, 再归约到显然成立的 $L=\Lambda$ 情形.

由 (2.1.13.1) 及通常的论证（参见 SGA~4, XVII, 5.2.11）可知, 若 $K\in\operatorname{Ob}D(X_\eta,\Lambda)$ 具有至多为 $d$ 的有限 Tor-维数, 其中 $d\in\mathbb Z$（SGA~4, XVII, 4.1.9）, 则 $R\Psi_\eta(K)$ 的 Tor-维数也至多为 $d$.
